\makeatletter
\def\input@path{{preamble/}{../preamble/}{../../preamble/}{../../../preamble/}{../../../../book/preamble/}}
\makeatother

\input{preamble_exams.tex}
% Add only the graph inputs used by the adapted key, if any.
\usepackage{varwidth}
\newcommand{\CreditDescription}[1]{\begin{varwidth}[t]{.72\linewidth}\raggedright#1\end{varwidth}}

\title{[Source exam title] --- Solutions}
\author{}
\date{}

\begin{document}
\maketitle

\setcounter{tocdepth}{1}
\phantomsection
\label{toc}
\tableofcontents

\section*{Learning objective}
[Source learning objective]

\section*{Criteria assessment}
[Source assessment instructions and grading policies]

% Use separate thresholds when the source distinguishes them.
[Source passing thresholds, if applicable]

\section*{Evaluated criteria}
\begin{description}
  \item[{[Criterion ID A]}] [Source criterion description]
  \item[{[Criterion ID B]}] [Source criterion description]
\end{description}

% Adapt one section per source question. Use the guide for explicit numbering.
% Single-task layout example: replace the entire example with the source task.
% Its algebra and point values are illustrative, not default scoring rules.
\sectionScored
  {Linear equation and substitution}
  {[Criterion ID A]-\sml{3}}
  {sec:question-single}

(i) Solve $2(x-3)=10$. \score{2}
(ii) Using your value of $x$, calculate $3x+1$. \score{1}

\subsection*{Solution \score{3}}
\label{sol:question-single}
\textbf{(i)} Expanding gives $2x-6=10$. Add 6 and divide by 2:
\[
2x=16,\qquad x=8.
\]
\textbf{(ii)} Substitute $x=8$: $3(8)+1=24+1=25$.

% Insert rubrics only in the final authoring pass after ALL solutions are
% complete and mathematically checked. Derive new rows from the actual work.
\par\smallskip
\noindent\begin{minipage}{\linewidth}\centering
\textbf{Proposed partial credit.}\par\smallskip
\begin{tabular}{@{}l r@{}}\toprule
\CreditDescription{\textbf{Assessed step or criterion}} & \textbf{Points}\\\midrule
\CreditDescription{(i): expanding to $2x-6=10$} & 0.5\\
\CreditDescription{(i): isolating $x$ by valid inverse operations to obtain $x=8$} & 1.5\\
\CreditDescription{(ii): substituting the obtained value of $x$ into $3x+1$} & 0.5\\
\CreditDescription{(ii): evaluating the resulting expression correctly} & 0.5\\
\midrule
\CreditDescription{\textbf{Total}} & \textbf{3}\\
\bottomrule\end{tabular}
\par\smallskip
Subtotals: (i) 2 points; (ii) 1 point.
\end{minipage}\par\smallskip

% Equivalent methods, such as dividing before expanding, earn corresponding
% credit; apply the governing marking policy to arithmetic-error follow-through.

\secInicio[][sec:question-single]

% Multipart pattern
\sectionScored
  {[Source question title]}
  {[Criterion ID A]-\sml{[Criterion A points]} \quad [Criterion ID B]-\sml{[Criterion B points]}}
  {sec:question-multi}

[Original shared context, data, and figures]

\QuestionsTasks[sec:question-multi]{
  {[Criterion A number]}/{sol:question-multi-first}/
    {[Source first subpart label and prompt] \score{[First subpart points]}},
  {[Criterion B number]}/{sol:question-multi-second}/
    {[Source second subpart label and prompt] \score{[Second subpart points]}}
}

\subsection*{Solution [Source first subpart label] \score{[First subpart points]}}
\label{sol:question-multi-first}
[Method, intermediate work, and interpreted result for the first subpart]

% Final authoring pass: insert this task's rubric here using the layout above.
\secInicio[][sec:question-multi]

\subsection*{Solution [Source second subpart label] \score{[Second subpart points]}}
\label{sol:question-multi-second}
[Method, intermediate work, and interpreted result for the second subpart]

% Final authoring pass: insert this task's rubric here using the layout above.
\secInicio[][sec:question-multi]
\end{document}
